\chapter{Trabalhos Relacionados}
\label{cap:trabalhos-relacionados}
%CORRIGIDO
Neste capítulo, são apresentados trabalhos que possuem similaridade com a proposta deste projeto, seja pelo domínio da aplicação, referindo-se ao setor veterinário, seja pela abordagem tecnológica, tratando-se de aplicações móveis ou de sincronização de dados.

\section{Prontuário Eletrônico Veterinário com Geolocalização (Pimazzoni, 2023)}
\label{sec:trabalho-relacionado-a}

%CORRIGIDO
O trabalho de \citeonline{pimazzoni2023} apresenta o desenvolvimento de um sistema de informações em saúde veterinária focado na integração de prontuários eletrônicos e geolocalização, com o objetivo de mapear áreas de maior incidência de zoonoses. A pesquisa destaca a necessidade de centralizar dados clínicos para melhorar a tomada de decisão do médico veterinário, além de estimular a atualização da base de dados com as infecções interespécie e seus respectivos vetores de transmissão, podendo tornar a prevenção e o tratamento futuro dessas doenças mais fáceis.

%CORRIGIDO
A aplicação VetWeb, desenvolvida a partir da pesquisa em questão, baseou-se no problema da escassez de dados públicos que tratem de doenças transmitidas entre animais e humanos, essas informações poderiam ser utilizadas posteriormente para pesquisas em saúde e disponibilizadas à população em geral, bem como no atraso gerado pela criação de prontuários manuais. O sistema possui telas voltadas para o cadastro de tutores, de pacientes e para a criação de prontuários clínicos, destinados à prescrição de exames, procedimentos e administração de vacinas. Além disso, cada paciente cadastrado na aplicação é devidamente localizado no mapa a partir do CEP fornecido em seu cadastro.

%CORRIGIDO
O sistema utiliza uma arquitetura cliente-servidor, além de estar disponível como aplicação \textit{mobile} e sistema \textit{web}, desenvolvido utilizando a linguagem PHP.




\section{Plataforma SimplesVet (Solução Comercial)}
\label{sec:trabalho-relacionado-b}
%CORRIGIDO
O SimplesVet \cite{simplesvet} é um dos principais \textit{softwares} de gestão veterinária no Brasil. Ele oferece módulos de prontuário, controle de estoque e financeiro, protocolos de vacina e planos pagos, operando majoritariamente no modelo \textit{Software as a Service} (SaaS). Essa aplicação auxilia no dia a dia dos veterinários, facilitando o trabalho em clínicas veterinárias, hospitais veterinários e \textit{pet shops}. É uma aplicação focada em empresas do ramo, pois além de fornecer as ferramentas para o uso do veterinário, ela também foca no controle financeiro da empresa que assina a plataforma, oferecendo telas para controle financeiro, como também estoque e serviços. 

%CORRIGIDO
Sobre a arquitetura do sistema, o SimplesVet foca em uma arquitetura baseada em nuvem, eliminando a necessidade de servidores locais. Além disso, os dados salvos no servidor da aplicação, vindos de diferentes pontos, são centralizados em uma única base de dados, como também a entrega desses dados em \textit{dashboards} profissionais, para medir o sucesso do SimplesVet no mercado, com o uso de métricas de desempenho.

%CORRIGIDO
Apesar de fornecer funcionalidades voltadas para o mercado veterinário, o SimplesVet opera no modelo SaaS, com os dados centralizados em nuvem, o que pressupõe a disponibilidade de conexão com a internet para o seu funcionamento. Dessa forma, a solução não é voltada ao público que atua em locais sem conectividade, como áreas isoladas de grandes cidades e regiões interioranas, cenário frequente no atendimento veterinário domiciliar. 



    

%CORRIGIDO
\section{Comparação entre os Trabalhos Relacionados}
\label{sec:trabalhos-relacionados-comparacao}
%CORRIGIDO
A Tabela~\ref{tab:trabalhos-relacionados} apresenta uma comparação entre os trabalhos relacionados e o RoadVet, considerando a adoção da arquitetura \textit{offline-first}, a presença de sincronização assíncrona e a portabilidade para dispositivos móveis.

%CORRIGIDO
\begin{table}[H]
    \centering
    \caption{Comparação entre os Trabalhos Relacionados}
    \label{tab:trabalhos-relacionados}
    \begin{tabular}{lccc}
        \toprule
        \textbf{Trabalho} & \textbf{Arquitetura \textit{offline-first}} & \textbf{Sincronização assíncrona} & \textbf{Portabilidade móvel} \\
        \midrule
        VetWeb     & Não & Não & Sim \\
        SimplesVet & Não & Não & Sim \\
        RoadVet    & Sim & Sim & Sim \\
        \bottomrule
    \end{tabular}
    \fonte{Elaboração própria.}
\end{table}


